\par\noindent\textbf{(1)} 惯性系速度平方为ṙ²+Ω²r²，L=m(ṙ²+Ω²r²)/2−k(r−r0)²/2，Euler方程 m r̈=mΩ²r−k(r−r0)−bṙ。角向约束转矩 τ=2mΩrṙ，保持Ω不变的电机功率 τΩ=2mΩ²rṙ，因此机械能本身一般不守恒。
\par\noindent\textbf{(2)} 静止平衡 r*=kr0/(k−mΩ²)，只有 k>mΩ² 才给有限正半径。有效势 V\_eff=k(r−r0)²/2−mΩ²r²/2，其二阶导数 k−mΩ²>0 恰是稳定条件；临界Ω\_c=√(k/m)时不再有正径向平衡。
\par\noindent\textbf{(3)} 设r=r*+ξ，线性扰动 mξ̈+bξ̇+(k−mΩ²)ξ=0，本征根 s\_±=[−b±√(b²−4m(k−mΩ²))]/(2m)。b=0时固有频率√(k/m−Ω²)，有阻尼时按判别式分欠阻尼与过阻尼；Ω超过临界使根有正实部，阻尼无法恢复线性稳定。
\par\noindent\textbf{(4)} 无阻尼时约化Jacobi量J=mṙ²/2+V\_eff(r)守恒：乘径向方程以ṙ得J̇=0。真实惯性系能量E=mṙ²/2+mΩ²r²/2+k(r−r0)²/2，与J差mΩ²r²，因此Ė=2mΩ²rṙ=τΩ，与外电机功精确对应。
